\chapter{Methodology}\label{cha:methodology}

This chapter defines the operators, projections, and diagnostics used throughout
the thesis, and states what each object measures and how it enters the
finite-sample and finite-width experiments. Detailed derivations are given in the
appendices.

\section{Mathematical setting}

\subsection{Domain and measure}

We identify the unit circle \(S^1\) with the angular domain \(\Omega=[0,2\pi)\), equipped
with the uniform probability measure
\begin{equation}
	d\mu(\varphi)=\frac{d\varphi}{2\pi}.
	\label{eq:method-uniform-measure}
\end{equation}
The embedding into \(\mathbb R^2\) is given by
\begin{equation}
	x(\varphi)=(\cos\varphi,\sin\varphi)\in S^1\subset\mathbb R^2.
	\label{eq:method-circle-embedding}
\end{equation}
This parametrization makes the rotational symmetry of the problem explicit and provides the
natural Fourier basis for \(L^2(\mu)\).

\subsection{Network model}

We consider a one-hidden-layer ReLU network in NTK parameterization. For inputs
\(x\in S^1\subset\mathbb{R}^2\),
\begin{equation}
	f_m(x;\theta)=\frac{1}{\sqrt m}\sum_{i=1}^m a_i \sigma(w_i^\top x+b_i),
	\label{eq:method-network-model}
\end{equation}
where \(\theta=\{(a_i,w_i,b_i)\}_{i=1}^m\), \(\sigma(t)=t_+=\max\{t,0\}\), and the parameters are
initialized according to
\begin{equation}
	a_i \sim \mathcal{N}(0,1),
	\qquad
	w_i \sim \mathcal{N}(0,I_2),
	\qquad
	b_i(0)=0,
	\label{eq:method-network-initialization}
\end{equation}
independently across \(i=1,\dots,m\). The bias parameters are initialized at zero but remain
trainable throughout.

\subsection{Least-squares loss and residual dynamics}
\label{subsec:least-squares-loss-and-residual-dynamics}

Let \(y\in L^2(\mu)\) be the target function, and write
\[
	f_t(\varphi):=f_m(x(\varphi);\theta(t))
\]
for the network output at training time \(t\). We consider the least-squares loss
\begin{equation}
	\mathcal{L}(\theta)
	=
	\frac{1}{2}\|f_m(\cdot;\theta)-y\|_{L^2(\mu)}^2
	=
	\frac{1}{2}\int_0^{2\pi}\bigl(f_m(\varphi;\theta)-y(\varphi)\bigr)^2\,d\mu(\varphi).
	\label{eq:method-loss-functional}
\end{equation}
We define the residual
\begin{equation}
	r_t(\varphi)=f_t(\varphi)-y(\varphi).
	\label{eq:method-residual-function}
\end{equation}

At finite width, the empirical neural tangent kernel at time \(t\) is
\begin{equation}
	\Theta_t^{(m)}(\varphi,\psi)
	=
	\bigl\langle \nabla_\theta f_m(\varphi;\theta(t)),
	\nabla_\theta f_m(\psi;\theta(t))\bigr\rangle,
	\label{eq:method-empirical-ntk}
\end{equation}
and the corresponding integral operator \(T_t^{(m)}:L^2(\mu)\to L^2(\mu)\) is
\begin{equation}
	(T_t^{(m)}f)(\varphi)
	=
	\int_0^{2\pi}\Theta_t^{(m)}(\varphi,\psi)f(\psi)\,d\mu(\psi).
	\label{eq:method-time-dependent-kernel-operator}
\end{equation}
Under gradient flow, the residual evolves according to
\begin{equation}
	\partial_t r_t=-T_t^{(m)}r_t.
	\label{eq:method-exact-residual-flow}
\end{equation}
Consequently, the loss evolves as
\begin{equation}
	\frac{d}{dt}\frac12\|r_t\|_{L^2(\mu)}^2
	=
	-\langle r_t,T_t^{(m)}r_t\rangle_{L^2(\mu)}.
	\label{eq:method-loss-decay-identity}
\end{equation}
Loss decay therefore depends not only on the eigenspaces of \(T_t^{(m)}\), but
also on how strongly the operator acts on the directions where the residual has
mass.

This equation contains both the frozen and evolving regimes. In the frozen regime,
\(T_t^{(m)}\) is replaced by its initial value
\begin{equation}
	T_0^{(m)}:=T_t^{(m)}\big|_{t=0},
	\label{eq:method-frozen-operator}
\end{equation}
whereas in the evolving regime the full time dependence of \(T_t^{(m)}\) is retained.

\section{Continuum reference model}

\subsection{Infinite-width tangent operator}

The first object studied in the thesis is the infinite-width limiting tangent operator
\begin{equation}
	T_\infty:=\lim_{m\to\infty} T_0^{(m)},
	\label{eq:method-infinite-width-operator}
\end{equation}
whose kernel is denoted by \(\Theta\). Under the initialization
\eqref{eq:method-network-initialization}, the limiting kernel exists and is deterministic.
The explicit formula for \(\Theta\) and the corresponding Fourier eigenvalues are stated in
Chapter~\ref{cha:results} and derived in Appendix~\ref{app:continuum-kernel}.

\subsection{Translation invariance and Fourier diagonalization}

Because the input distribution is uniform and the limiting kernel is rotation-invariant,
\(\Theta(\varphi,\psi)\) depends only on the angular difference. The limiting operator
therefore takes the convolution form
\begin{equation}
	(T_\infty f)(\varphi)
	=
	\int_0^{2\pi}\Theta(\varphi-\psi)f(\psi)\,d\mu(\psi),
	\label{eq:method-continuum-operator}
\end{equation}
and is diagonalized by the real Fourier basis of \(L^2(\mu)\). Its eigenvalues determine the
decay rates of the Fourier components of the residual in the ideal fixed-kernel continuum
model.

This continuum Fourier decomposition serves as the reference point for the finite-sample and
finite-width analyses below.

\subsection{Real Fourier basis and truncated subspaces}
\label{subsec:real-fourier-basis-truncates-subspaces}

We define the real Fourier basis of \(L^2(\mu)\) by
\begin{equation}
	\phi_0(\varphi)=1,
	\label{eq:method-fourier-basis-phi0}
\end{equation}
and, for each \(k\geq 1\),
\begin{equation}
	\phi_{k,c}(\varphi)=\sqrt{2}\cos(k\varphi),
	\qquad
	\phi_{k,s}(\varphi)=\sqrt{2}\sin(k\varphi).
	\label{eq:method-fourier-basis-modes}
\end{equation}
For a fixed cutoff \(K\geq 1\), we consider the truncated subspace
\begin{equation}
	\mathcal{H}_K=\operatorname{span}\{\phi_0,\phi_{k,c},\phi_{k,s}:1\leq k\leq K\},
	\label{eq:method-low-frequency-subspace}
\end{equation}
with dimension
\begin{equation}
	d=2K+1.
	\label{eq:method-block-dimension}
\end{equation}
Let
\begin{equation}
	P_K:L^2(\mu)\to \mathcal{H}_K
	\label{eq:method-projection-PK}
\end{equation}
denote the orthogonal projection onto $\mathcal H_K$.

The truncation to \(\mathcal{H}_K\) is used throughout the thesis to reduce the continuum problem to a
finite-dimensional low-frequency block that can be compared across the continuum, sampled,
frozen finite-width, and evolving finite-width settings.

\section{Finite-sample methodology}
\label{sec:finite-sample-methodology}

\subsection{Finite-sample question}

We first study the effect of replacing the measure \(\mu\) by finitely many sampled points. Let
\begin{equation}
	\varphi_1,\dots,\varphi_n \overset{\mathrm{iid}}{\sim} \mathrm{Unif}[0,2\pi).
	\label{eq:method-sampled-angles}
\end{equation}
The main question is whether the diagonal Fourier decomposition of the continuum operator
remains approximately valid after this replacement.

\subsection{Sampled objects}

To study the sampled operator on \(\mathcal{H}_K\), we introduce the feature matrix
\(\Phi\in\mathbb{R}^{n\times d}\),
\begin{equation}
	\Phi_{i,p}=\phi_p(\varphi_i),
	\qquad
	1\leq i\leq n,\quad 1\leq p\leq d,
	\label{eq:method-feature-matrix}
\end{equation}
and the normalized sampled kernel matrix \(A\) with
\begin{equation}
	A_{ij}=\frac{1}{n}\Theta(\varphi_i-\varphi_j).
	\label{eq:method-empirical-operator}
\end{equation}
From these we form the Gram matrix
\begin{equation}
	G=\frac{1}{n}\Phi^\top\Phi,
	\label{eq:method-gram-matrix}
\end{equation}
and the compressed operator matrix
\begin{equation}
	H=\frac{1}{n}\Phi^\top A\Phi.
	\label{eq:method-compressed-operator}
\end{equation}
We also introduce the diagonal matrix
\begin{equation}
	\Lambda^{(n)}=\operatorname{diag}(\lambda_p^{(n)}),
	\label{eq:method-finite-sample-eigenvalue-matrix}
\end{equation}
whose entries are the finite-sample corrected Fourier eigenvalues defined in
Chapter~\ref{cha:results}.

\subsection{Quantities to be controlled}
The finite-sample analysis controls how far three sampled quantities deviate from
the values they would take if the sampled Fourier structure were exact. Each is
written as the difference between the sampled quantity and its ideal value:
\begin{equation}
	\underbrace{G-I_d}_{\text{basis vs.\ orthonormal}},
	\qquad
	\underbrace{H-\Lambda^{(n)}}_{\text{operator vs.\ diagonal}},
	\qquad
	\underbrace{A\Phi-\Phi\Lambda^{(n)}}_{\text{modes vs.\ eigenvectors}}.
	\label{eq:method-finite-sample-objects}
\end{equation}
The first is small when the sampled Fourier basis is close to orthonormal, so its
Gram matrix \(G\) is close to the identity. The second is small when the sampled
operator restricted to \(\mathcal H_K\), represented by \(H\), is close to the
diagonal finite-sample prediction \(\Lambda^{(n)}\). The third is small when the
sampled Fourier modes are close to being eigenvectors of the sampled operator,
i.e.\ applying the operator \(A\) to the modes \(\Phi\) reproduces them scaled by
the predicted eigenvalues. Together these are the quantities needed to answer Q1.

\section{Fourier-block preconditioning methodology}
\label{sec:method-preconditioning}

The finite-sample analysis gives a diagonal reference for the sampled operator on
a fixed Fourier block. We use it to define a preconditioning experiment that tests
whether the predicted Fourier eigenvalues explain the observed rate separation
between retained Fourier modes.

Fix a retained Fourier block
\[
	\mathcal H_K
	=
	\operatorname{span}\{\phi_0,\phi_{1,c},\phi_{1,s},\dots,\phi_{K,c},\phi_{K,s}\},
\]
with dimension \(d=2K+1\). Let
\[
	\Phi\in\mathbb R^{n\times d}
\]
be the sampled Fourier feature matrix, and let \(A\in\mathbb R^{n\times n}\) be
the normalized sampled kernel operator. Recall that
\begin{equation}
	G=\frac1n\Phi^\top\Phi,
	\qquad
	H=\frac1n\Phi^\top A\Phi.
	\label{eq:method-preconditioning-GH}
\end{equation}
The matrix \(G\) is the sampled Fourier Gram matrix, while \(H\) records the
action of the sampled operator \(A\) on the retained Fourier basis.
When \(G\) is invertible, define
\begin{equation}
	C_n^{(K)}
	:=
	G^{-1}H.
	\label{eq:method-sampled-coefficient-operator}
\end{equation}
This matrix describes what the sampled operator \(A\) does to the Fourier
coefficients on the retained block. Indeed, if \(u=\Phi z\), then \(z\) is the
coefficient vector of \(u\) in the sampled Fourier basis. After applying \(A\),
the least-squares coefficients of \(Au\) in the same sampled Fourier basis are
\begin{equation}
	G^{-1}\frac1n\Phi^\top (Au)
	=
	G^{-1}\frac1n\Phi^\top A\Phi z
	=
	C_n^{(K)}z.
	\label{eq:method-sampled-coefficient-action}
\end{equation}
Thus \(C_n^{(K)}\) is the object that should be compared with the diagonal
finite-sample prediction
\begin{equation}
	\Lambda^{(n)}
	=
	\operatorname{diag}(\lambda_p^{(n)}).
	\label{eq:method-preconditioning-lambda-n}
\end{equation}
The finite-sample theory predicts that, on the retained block,
\begin{equation}
	C_n^{(K)}
	\approx
	\Lambda^{(n)}.
	\label{eq:method-coefficient-operator-approx}
\end{equation}
This means that, on the retained Fourier block, the sampled operator is expected
to act mainly by rescaling each Fourier coefficient by its predicted eigenvalue,
with only small mixing between different retained Fourier modes.

We now define two preconditioners from this comparison. First note that, for any
sampled vector \(v\in\mathbb R^n\),
\[
	G^{-1}\frac1n\Phi^\top v
\]
gives the least-squares Fourier coefficients of the projection of \(v\) onto
the sampled span of the retained modes. Thus, to correct the action of \(A\) on
the retained block, we first apply \(A\), then project back onto the retained
Fourier span, and finally rescale the resulting Fourier coefficients.

The theory preconditioner uses the diagonal prediction \(\Lambda^{(n)}\):
\begin{equation}
	P_{\mathrm{th}}
	:=
	\frac1n
	\Phi
	\bigl(\Lambda^{(n)}\bigr)^{-1}
	G^{-1}
	\Phi^\top .
	\label{eq:method-theory-preconditioner}
\end{equation}
This operator projects a sampled vector onto the retained Fourier span and
rescales each retained Fourier coefficient by the inverse predicted eigenvalue.
It uses the analytic prediction \(\Lambda^{(n)}\), and therefore does not use
the observed matrix \(C_n^{(K)}\).

The empirical preconditioner uses the observed retained block:
\begin{equation}
	P_{\mathrm{emp}}
	:=
	\frac1n
	\Phi
	\bigl(C_n^{(K)}\bigr)^{-1}
	G^{-1}
	\Phi^\top .
	\label{eq:method-empirical-preconditioner}
\end{equation}
This operator uses the actual sampled matrix \(C_n^{(K)}\). It therefore
corrects both the diagonal scaling and the mixing between retained Fourier modes
that is present for the particular sampled points.

We compare three sampled dynamics:
\begin{equation}
	B_{\mathrm{base}}:=A,
	\qquad
	B_{\mathrm{th}}:=P_{\mathrm{th}}A,
	\qquad
	B_{\mathrm{emp}}:=P_{\mathrm{emp}}A.
	\label{eq:method-preconditioned-operators}
\end{equation}
The baseline \(B_{\mathrm{base}}\) is the original sampled kernel operator. The
operators \(B_{\mathrm{th}}\) and \(B_{\mathrm{emp}}\) apply the sampled kernel
operator and then correct its action through the retained Fourier block.

For sampled predictions \(f_t\in\mathbb R^n\), targets \(y\in\mathbb R^n\), and
residuals
\[
	r_t=f_t-y,
\]
we compare the gradient-flow dynamics
\begin{equation}
	\frac{d r_t}{dt}
	=
	- B r_t,
	\label{eq:method-sampled-operator-flow}
\end{equation}
with
\[
	B\in\{B_{\mathrm{base}},B_{\mathrm{th}},B_{\mathrm{emp}}\}.
\]

The effect is easiest to see on a vector \(u=\Phi z\) in the retained Fourier
span. Applying \(A\) and projecting back to the retained block gives the
coefficient vector \(C_n^{(K)}z\). Therefore
\begin{align}
	B_{\mathrm{th}}\Phi z
	 & =
	\Phi
	\bigl(\Lambda^{(n)}\bigr)^{-1}
	C_n^{(K)}z,
	\label{eq:method-theory-preconditioned-block-action}
\end{align}
while
\begin{align}
	B_{\mathrm{emp}}\Phi z
	 & =
	\Phi
	\bigl(C_n^{(K)}\bigr)^{-1}
	C_n^{(K)}z
	=
	\Phi z.
	\label{eq:method-empirical-preconditioned-block-action}
\end{align}
Thus on the retained block the empirical preconditioner reduces the action to the
identity: every retained Fourier mode is driven at the same rate, so the
frequency-dependent rate separation is removed, provided \(C_n^{(K)}\) is
invertible. The theory preconditioner achieves this only insofar as the observed
block \(C_n^{(K)}\) is close to the diagonal prediction \(\Lambda^{(n)}\); from
\eqref{eq:method-theory-preconditioned-block-action}, it leaves the residual
factor \((\Lambda^{(n)})^{-1}C_n^{(K)}\), which equals the identity exactly when
\(C_n^{(K)}=\Lambda^{(n)}\).

In the numerical experiments, ill-conditioned inverses are replaced by small
ridge-regularized inverses when needed, for example \((\Lambda^{(n)}+\tau I)^{-1}\)
or \((C_n^{(K)}+\tau I)^{-1}\) with \(\tau>0\). If \(B_{\mathrm{th}}\) and
\(B_{\mathrm{emp}}\) behave similarly, the predicted Fourier eigenvalues explain
most of the rate differences on the retained block; if \(B_{\mathrm{emp}}\) is
noticeably better, the sampled points introduce mixing between retained Fourier
modes that the diagonal prediction does not capture.

\section{Finite-width frozen-kernel methodology}

\subsection{Finite-width question}

We next remove sampling randomness and return to the continuum input space \(S^1\), but keep
the network width finite. The question is then whether the low-frequency block of the frozen
finite-width tangent operator remains close to the continuum Fourier block predicted by the
infinite-width theory.

\subsection{One-neuron tangent features and operators}

For each neuron \(r\), define its tangent feature at initialization by
\begin{equation}
	\psi_r(\varphi)
	=
	\nabla_{(a_r,w_r,b_r)}
	\bigl(a_r\sigma(w_r^\top x(\varphi)+b_r)\bigr).
	\label{eq:method-one-neuron-tangent-feature}
\end{equation}
The corresponding one-neuron kernel contribution is
\begin{equation}
	\Theta^{(r)}(\varphi,\psi)
	=
	\psi_r(\varphi)^\top \psi_r(\psi),
	\label{eq:method-one-neuron-kernel}
\end{equation}
and the associated one-neuron operator is
\begin{equation}
	(T^{(r)}f)(\varphi)
	=
	\int_0^{2\pi}\Theta^{(r)}(\varphi,\psi)f(\psi)\,d\mu(\psi).
	\label{eq:method-one-neuron-operator}
\end{equation}
At initialization, the frozen finite-width tangent operator is the average
\begin{equation}
	T_0^{(m)}=\frac1m\sum_{r=1}^m T^{(r)}.
	\label{eq:method-finite-width-frozen-average}
\end{equation}

\subsection{Projected frozen operator}

Let \( P_K:L^2(\mu)\to \mathcal H_K \)
denote the orthogonal projection onto the truncated Fourier subspace
\(\mathcal H_K\). The low-frequency block of the frozen finite-width operator is
\begin{equation}
	B_m^{(K)}:=P_KT_0^{(m)}P_K.
	\label{eq:method-finite-width-block}
\end{equation}
The corresponding continuum reference block is
\begin{equation}
	\Lambda_K:=P_KT_\infty P_K
	=
	\operatorname{diag}(\lambda_0,\lambda_1,\lambda_1,\dots,\lambda_K,\lambda_K).
	\label{eq:method-continuum-fourier-block}
\end{equation}
The finite-width frozen analysis asks whether \(B_m^{(K)}\) remains close to \(\Lambda_K\)
with high probability as \(m\) increases.

\subsection{Matrix representation of the frozen block}

Using the ordered basis
\[
	(\phi_0,\phi_{1,c},\phi_{1,s},\dots,\phi_{K,c},\phi_{K,s}),
\]
the projected operator \(B_m^{(K)}\) is represented by the matrix with entries
\begin{equation}
	(B_m^{(K)})_{pq}
	=
	\langle \phi_p,T_0^{(m)}\phi_q\rangle_{L^2(\mu)}.
	\label{eq:method-finite-width-block-entry}
\end{equation}
Equivalently,
\begin{equation}
	(B_m^{(K)})_{pq}
	=
	\frac1m\sum_{r=1}^m \xi_{r,pq},
	\qquad
	\xi_{r,pq}:=\langle \phi_p,T^{(r)}\phi_q\rangle,
	\label{eq:method-finite-width-entry-average}
\end{equation}
so that \(B_m^{(K)}\) is an empirical average of \(m\) i.i.d.\ one-neuron
matrices. Because the entries are averages of independent terms, they concentrate
around their mean as \(m\) grows, which is what the finite-width frozen
concentration result in the next chapter exploits.

\subsection{Fourier-plane alignment diagnostics}
\label{subsec:method-projector-error}

The continuum operator is diagonal in the Fourier basis. For each frequency
\(k\ge 1\), the cosine and sine modes have the same continuum eigenvalue
\(\lambda_k\). Because \(\phi_{k,c}\) and \(\phi_{k,s}\) share the eigenvalue
\(\lambda_k\), the continuum operator does not single out either one; only the
plane they span is determined. The object to compare against is therefore this
two-dimensional frequency subspace,
\[
	\mathcal F_k
	:=
	\operatorname{span}\{\phi_{k,c},\phi_{k,s}\},
	\qquad k\ge 1,
\]
For the constant mode, we set
\[
	\mathcal F_0:=\operatorname{span}\{\phi_0\}.
\]
Thus, the dimensionality \(d_0=1\) and \(d_k=2\) for \(k\ge 1\).

We use subspace-alignment diagnostics to compare these reference Fourier
subspaces with matched eigenspaces of the finite-width tangent operator. In the
frozen case this operator is \(T_0^{(m)}\). In the evolving case it is the
time-dependent operator \(T_t^{(m)}\).

The operator is defined on the continuum \(S^1\), but it cannot be eigen-decomposed
in closed form at finite width, so for the numerical diagnostics we discretise it
onto a uniform grid of \(N\) points,
\[
	\theta_j=\frac{2\pi j}{N},\qquad j=0,\dots,N-1,
\]
and carry out the eigendecomposition and subspace comparisons on this grid.

All vectors and matrices in this diagnostic are therefore finite-dimensional
representations of continuum functions on this grid.

Let \(U_k\in\mathbb R^{N\times d_k}\) be an orthonormal basis for the grid
representation of \(\mathcal F_k\), and define the corresponding orthogonal
projector
\begin{equation}
	P_k := U_kU_k^\top .
	\label{eq:method-fourier-projector}
\end{equation}

At each training time \(t\), we eigendecompose the grid representation of the
tangent operator, an \(N\times N\) symmetric matrix, and write its orthonormal
eigenvectors as \(v_1(t),\dots,v_N(t)\). For each eigenvector \(v_j(t)\), the
quantity
\[
	\|U_k^\top v_j(t)\|_2^2 \in [0,1]
\]
is the squared length of its projection onto \(\mathcal F_k\): it is \(1\) when
\(v_j(t)\) lies entirely in \(\mathcal F_k\) and \(0\) when it is orthogonal to
it. We build the matched empirical subspace \(\widehat{\mathcal F}_k(t)\) by taking
the \(d_k\) eigenvectors with the largest such values, the ones lying closest to
\(\mathcal F_k\). If \(\widehat U_k(t)\in\mathbb R^{N\times d_k}\) is an
orthonormal basis of \(\widehat{\mathcal F}_k(t)\), define
\begin{equation}
	\widehat P_k(t)
	:=
	\widehat U_k(t)\widehat U_k(t)^\top .
	\label{eq:method-empirical-projector}
\end{equation}

The projection error is
\begin{equation}
	\varepsilon_{k,F}^{(m)}(t)
	:=
	\|\widehat P_k(t)-P_k\|_F .
	\label{eq:method-projector-errors}
\end{equation}
The projectors \(P_k\) and \(\widehat P_k(t)\) depend only on the subspaces, not
on the bases chosen for them, so this distance is invariant to sign changes of
eigenvectors and to rotations of the cosine-sine basis inside a two-dimensional
Fourier plane, neither of which should count as a mismatch. It therefore measures
the geometric distance between the matched eigenspace and the reference Fourier
subspace, and \eqref{eq:method-projector-angle-identity} below expresses it
through the principal angles. For the frozen operator, we write
\[
	\varepsilon_{k,F}^{(m)}:=\varepsilon_{k,F}^{(m)}(0).
\]

We also report principal angles between
\(\widehat{\mathcal F}_k(t)\) and \(\mathcal F_k\). Principal angles are a
standard way to compare finite-dimensional subspaces
\citep{bjorck1973numerical}. If \(U_k\) and \(\widehat U_k(t)\) are orthonormal
bases for the two subspaces, then the cosines of the principal angles are the
singular values of \(U_k^\top \widehat U_k(t)\). We define
\[
	0\le \theta_{k,1}(t)\le \cdots \le \theta_{k,d_k}(t)\le \frac{\pi}{2}
\]
by
\begin{equation}
	\cos\theta_{k,i}(t)
	=
	s_i(U_k^\top \widehat U_k(t)),
	\qquad i=1,\dots,d_k,
	\label{eq:method-principal-angles}
\end{equation}
where \(s_i(\cdot)\) is the \(i\)-th largest singular value, so the cosines of the
principal angles are the singular values of \(U_k^\top\widehat U_k(t)\) ordered
from largest to smallest.

For \(k=0\), there is one principal angle because both spaces are
one-dimensional. For \(k\ge 1\), there are two principal angles because both
spaces are two-dimensional. The first angle measures the best aligned direction
between the two planes. The second angle measures the remaining independent
direction. Both angles must be small for the full Fourier plane to be well
aligned with the matched empirical eigenspace.

For the evolving-kernel experiments, we also summarize the principal angles by
the subspace cosine similarity
\begin{equation}
	s_k^{\mathrm{cos}}(t)
	:=
	\left(
	\frac{1}{d_k}
	\sum_{i=1}^{d_k}\cos^2\theta_{k,i}(t)
	\right)^{1/2}.
	\label{eq:method-subspace-cosine-similarity}
\end{equation}
This quantity is the root-mean-square of the cosines of the principal angles.
We use this aggregation to summarize one-dimensional and two-dimensional Fourier
subspaces with a single score per frequency. The score lies between \(0\) and
\(1\), with larger values indicating better alignment. It equals \(1\) exactly
when the matched empirical subspace is the same as the reference Fourier
subspace.

The projection distance and the principal angles are related by the standard
projector identity \citep{bjorck1973numerical}
\begin{equation}
	\|\widehat P_k(t)-P_k\|_F^2
	=
	2\sum_{i=1}^{d_k}\sin^2\theta_{k,i}(t).
	\label{eq:method-projector-angle-identity}
\end{equation}
Thus the projection error gives one scalar measure of subspace mismatch, while
the principal angles and subspace cosine similarity give more detailed views of
the same alignment.

\section{Finite-width evolving-kernel methodology}
\label{sec:finite-width-evolving-kernel-methodology}
\subsection{Projected evolving operator}

To study the actual training dynamics, we retain the time dependence of the
empirical tangent operator \(T_t^{(m)}\). Its block on the same retained
subspace \(\mathcal H_K\), the frequencies up to \(K\) used throughout, is
\begin{equation}
	C_t^{(K)}:=P_KT_t^{(m)}P_K.
	\label{eq:method-evolving-block}
\end{equation}

In the ordered Fourier basis of \(\mathcal H_K\), this block has entries
\begin{equation}
	(C_t^{(K)})_{pq}
	=
	\langle \phi_p,T_t^{(m)}\phi_q\rangle_{L^2(\mu)}.
	\label{eq:method-evolving-block-entry}
\end{equation}
At initialization,
\[
	C_0^{(K)}=B_m^{(K)}.
\]
If \(b_t=(I-P_K)r_t\) is the residual component outside the retained Fourier
block, then \(T_t^{(m)}b_t\) need not remain outside that block. Its projection
back into \(\mathcal H_K\) is
\begin{equation}
	L_t^{(K)}b_t
	=
	P_KT_t^{(m)}(I-P_K)r_t.
\end{equation}
We therefore define
\begin{equation}
	L_t^{(K)}:=P_KT_t^{(m)}(I-P_K).
	\label{eq:method-leakage-operator}
\end{equation}
This operator represents high-to-low frequency coupling. It measures how
residual components outside the retained block can influence the evolution of
the retained Fourier coefficients after the tangent operator is applied.

\subsection{Projected residual dynamics}

Let
\begin{equation}
	a_t:=P_Kr_t,
	\qquad
	b_t:=(I-P_K)r_t.
	\label{eq:method-low-high-residual-split}
\end{equation}
Projecting the exact residual equation \eqref{eq:method-exact-residual-flow} gives
\begin{equation}
	\partial_t a_t
	=
	-C_t^{(K)}a_t
	-
	L_t^{(K)}b_t.
	\label{eq:method-projected-residual-flow}
\end{equation}
Thus the low-frequency dynamics is driven by two terms: the operator restricted
to \(\mathcal H_K\), and the contribution from residual components outside
\(\mathcal H_K\).

\subsection{Frozen versus evolving comparison}

The frozen and evolving regimes differ only through the time dependence of the operator. In
the frozen regime,
\[
	C_t^{(K)}\equiv C_0^{(K)}=B_m^{(K)},
\]
whereas in the evolving regime the operator changes during training. This motivates the
decomposition
\begin{equation}
	C_t^{(K)}
	=
	\Lambda_K
	+
	\bigl(C_0^{(K)}-\Lambda_K\bigr)
	+
	\bigl(C_t^{(K)}-C_0^{(K)}\bigr),
	\label{eq:method-operator-decomposition}
\end{equation}
which separates three effects:

\begin{enumerate}
	\item the ideal infinite-width Fourier block \(\Lambda_K\),
	\item the finite-width initialization error \(C_0^{(K)}-\Lambda_K\),
	\item the time-dependent kernel drift \(C_t^{(K)}-C_0^{(K)}\).
\end{enumerate}

The frozen-kernel theory studies the second term. The evolving-kernel
experiments study the third term empirically: how much the tangent operator
changes during training, and how these changes affect the Fourier-block
structure.

\subsection{Frequency-wise spectral strength}
\label{subsec:method-spectral-strength}

The subspace-alignment diagnostics measure whether the eigenspaces of the
finite-width tangent operator remain close to the Fourier subspaces. They do
not, however, measure how strongly the operator acts on those subspaces.
For the residual dynamics, this distinction matters because the loss decay is
controlled by \(\langle r_t,T_t^{(m)}r_t\rangle_{L^2(\mu)}\), as shown in
\eqref{eq:method-loss-decay-identity}. Thus the loss decay depends not only on
the orientation of the eigenspaces, but also on the magnitude of the operator in
directions where the residual has mass.
For each Fourier frequency subspace \(\mathcal F_k\), with orthogonal projector
\(P_k\) and dimension \(d_k\), we define the average spectral strength
\begin{equation}
	\mu_{\mathcal F_k}(t)
	:=
	\frac{1}{d_k}
	\operatorname{tr}\!\left(P_k T_t^{(m)} P_k\right).
	\label{eq:method-frequency-spectral-strength}
\end{equation}
Equivalently, if \(U_k\) is an orthonormal basis for \(\mathcal F_k\), then
\begin{equation}
	\mu_{\mathcal F_k}(t)
	=
	\frac{1}{d_k}
	\operatorname{tr}\!\left(U_k^\top T_t^{(m)} U_k\right).
	\label{eq:method-frequency-spectral-strength-basis}
\end{equation}
This quantity is basis-invariant within \(\mathcal F_k\). If \(\mathcal F_k\)
is an invariant eigenspace of the operator, then
\(\mu_{\mathcal F_k}(t)\) equals the corresponding eigenvalue. In general, it
should instead be interpreted as the average kernel strength inside that
Fourier block.

For \(k\le K\), the same quantity can be read directly from the low-frequency
block \(C_t^{(K)}\). It is the average trace of the diagonal block corresponding
to \(\mathcal F_k\):
\begin{equation}
	\mu_{\mathcal F_k}(t)
	=
	\frac{1}{d_k}
	\operatorname{tr}\!\left((C_t^{(K)})_{\mathcal F_k,\mathcal F_k}\right).
	\label{eq:method-frequency-spectral-strength-block}
\end{equation}

We also use the cumulative low-frequency strength
\begin{equation}
	\mu_{\le K}(t)
	:=
	\frac{1}{2K+1}
	\operatorname{tr}\!\left(P_KT_t^{(m)}P_K\right)
	=
	\frac{1}{2K+1}
	\sum_{k=0}^K d_k\,\mu_{\mathcal F_k}(t).
	\label{eq:method-cumulative-spectral-strength}
\end{equation}
This summarizes the average tangent-kernel strength over the retained
low-frequency block.

\section{Summary of the methodology}

The methodology consists of three comparisons.

First, the continuum infinite-width operator \(T_\infty\) gives the ideal
Fourier-diagonal reference model on \(S^1\). In this model, the Fourier modes
are exact eigenfunctions and the eigenvalues determine the residual decay rates.

Second, the finite-sample analysis studies what changes when the continuum
measure is replaced by finitely many sampled points. The main objects are the
sampled Fourier Gram matrix \(G\), the matrix \(H\) describing the action of
the sampled operator on the retained Fourier basis, and the difference
\(A\Phi-\Phi\Lambda^{(n)}\) between the sampled operator action and the
diagonal Fourier prediction. These objects are used to answer Q1. The rate
separation between modes comes from the spread of eigenvalues, so rescaling
each mode by its predicted eigenvalue should cancel it; the preconditioning
experiment uses this to test how much of the separation the eigenvalues
explain.

Third, the finite-width analysis studies what changes when the infinite-width
operator is replaced by the tangent operator of a finite-width network. The
frozen analysis focuses on \(B_m^{(K)}=P_KT_0^{(m)}P_K\), the low-frequency block
of the tangent operator at initialization. The evolving analysis tracks the
time-dependent block \(C_t^{(K)}=P_KT_t^{(m)}P_K\) and the high-to-low coupling
\(L_t^{(K)}=P_KT_t^{(m)}(I-P_K)\).

For the evolving experiments, we use two complementary kinds of diagnostics.
Projection error, principal angles, and subspace cosine similarity measure
whether the eigenspaces of the tangent operator remain close to the Fourier
subspaces.
Frequency-wise spectral strength measures how strongly the tangent operator
acts on each Fourier block. Together, these diagnostics separate changes in
Fourier geometry from changes in the magnitude of the kernel on the retained
frequency blocks.
